\documentclass[10pt,a4paper]{amsart}
\usepackage[margin=19mm]{geometry}
\usepackage{amsmath,amssymb,mathtools}
\usepackage[scr=rsfs,cal=euler]{mathalfa}
\usepackage{xfrac}
\usepackage[unicode,hidelinks,bookmarks=false]{hyperref}
\newcommand{\ie}{i.e.\ }
\newcommand{\cf}{cf.\ }
\newcommand{\resp}{respectively\ }
\newcommand{\Subsection}{\subsection}
\providecommand{\nobreakdash}{\nobreak\hbox{-}}
\setlength{\emergencystretch}{2em}
\multlinegap0pt
\usepackage{bm}
\usepackage{bbm}
\usepackage{dsfont}
\usepackage{xspace}
\usepackage{cancel}
\usepackage{mathtools}
\usepackage{amsthm}
\makeatletter
\@ifundefined{theorem}{\newtheorem{theorem}{Theorem}}{}
\@ifundefined{lemma}{\newtheorem{lemma}[theorem]{Lemma}}{}
\makeatother
\title[Solutions with one dimensional concentration for a two dimen]{Solutions with one dimensional concentration for a two dimensional Gross-Pitaevskii model with general potential}
\author{Lipeng Duan, Suting Wei, Jun Yang}
\date{Source: arXiv:2602.20554v1 (CC BY 4.0)}
\begin{document}
\maketitle

\noindent\emph{Source and licence.} Solutions with one dimensional concentration for a two dimensional Gross-Pitaevskii model with general potential. Version arXiv:2602.20554v1 (CC BY 4.0), distributed under the Creative Commons Attribution 4.0 International licence, as recorded in the arXiv licence metadata for this exact version and shown on the abstract page https://arxiv.org/abs/2602.20554v1. Full rights notice: licenses/arxiv-2602.20554-CC-BY-4.0.txt.\par\smallskip

\noindent\textit{Editorial scope.} Counted unit: the Lemma whose source label is \texttt{linear lemma6-1} in the article (source0.tex lines 5114--5132, source label \texttt{linear lemma6-1}), delivered together with the article's own complete original proof of it (source0.tex lines 5135--5265, one \texttt{proof} environment of 131 lines). No mathematics is written, repaired or re-derived: statement and proof are the article's own text line for line. Required context retained uncounted: lambda* (lines 869--874); beta-bound (lines 1397--1404); linearphi1 (lines 5102--5105); linearphi2 (lines 5106--5109); linearphi3 (lines 5110--5112). Every definition, display and label of those lines is retained unchanged. Omitted from this package and not redistributed: 1--868, 875--1396, 1405--5101, 5113--5113, 5133--5134, 5266--9220. Nothing inside a retained excerpt is omitted or altered: every retained body is delivered exactly as the article's lines are written, so no omission receipt of any kind accompanies this package. Citations: the retained excerpts carry no citation command at all, so no bibliography entry of the article is transcribed and no reference is redistributed. Reference adaptation: none was needed and none was made; every reference inside the delivered unit resolves inside it, so no unresolved reference or citation is produced. Printed slips kept literal: no printed word, symbol, hypothesis or number of the counted statement or of its proof was altered. Layout and class: the article's own class is \texttt{amsart} with packages including amssymb, latexsym, amsmath, amsthm, graphicx, hyperref, titletoc, pdfsync, color, xcolor, mathrsfs, multirow, caption, appendix; the wrapper is the lane's pinned \texttt{amsart} editorial wrapper at 10pt on A4, which restates the theorem-like environments the retained text needs and adds the article's own macro definitions that the retained text needs (none), with extra wrapper packages (bm, bbm, dsfont, xspace, cancel, mathtools). The delivered numbering is the unit's own. Assets: no retained span contains a figure environment, a graphics-inclusion command, a PDF-page inclusion, a table or a TikZ picture; no image file is copied into this package, the package declares no asset dependency and no image file is redistributed with it. Licence: version arXiv:2602.20554v1 of this article is distributed under the Creative Commons Attribution 4.0 International licence, as recorded in the arXiv licence metadata for this exact version and shown on the abstract page https://arxiv.org/abs/2602.20554v1. A full rights notice is delivered at licenses/arxiv-2602.20554-CC-BY-4.0.txt. Characters: every retained excerpt is pure ASCII, so the delivered engine, XeLaTeX with its default Latin Modern text font, reports no missing character.
\begin{align}
\lambda_* = \frac{\lambda_0} {4\pi^2}
\qquad \text{and}\qquad
{\tilde\ell} = \int_0^{\ell} \big[1+ \epsilon^2 W(0, \theta)\big]^{\frac{1}{2}} {\mathrm d} \theta.
\label{lambda*}
\end{align}

\begin{align}
|\beta'|\,\leq\, C\beta,
\qquad
|\beta''|\,\leq\, C\beta,
\qquad
|\beta'''|\,\leq\, C\beta.
\label{beta-bound}
\end{align}

\begin{align} \label{linearphi1}
 \frac{\epsilon^2 4\pi^2}{{\tilde\ell}^2} \phi_{\tau\tau} + \phi_{xx} - \phi+3U^2 \phi = h,
\quad \text{for}\,(x, \tau)\in \mathfrak S,
\end{align}

\begin{align}
 \phi(x, 0) = \phi(x, 2\pi), \quad \phi_\tau(x, 0) = \phi_\tau(x, 2\pi), \quad \text{for}~x\in {\mathbb R},
\label{linearphi2}
\end{align}

\begin{align} \label{linearphi3}
 \int_{\mathbb R}\phi(x, \tau)\,\mathcal Z(x) \,{\mathrm d}x = \int_{\mathbb R}\phi(x, \tau)\,U'(x)\,{\mathrm d}x \,=\, 0, \quad \text{for}~\tau\in [0, 2\pi),
\end{align}

\begin{lemma}\label{linear lemma6-1}
 Suppose $\phi$ solves \eqref{linearphi1}-\eqref{linearphi3}. Then there exists a constant $C>0$, independent of $\epsilon$, such that
\begin{align} \label{norm h2 mathfrak S estimate}
 \|\phi\|_{H^2_*(\mathfrak S)} \leq C \|h\|_{L^2(\mathfrak S)},
\end{align}
where the  norm  $\|\cdot\|_{_{H^2_*(\mathfrak S)}}$ is given by
\begin{align} \label{norm h2 mathfrak S}
\|\phi\|_{H^2_*(\mathfrak S)}
= \left({\displaystyle\int}_{\mathfrak S} \Bigg[\, |\phi|^2+ |\phi_{x}|^2+ |\phi_{xx}|^2 +\Big|\frac{\epsilon}{\tilde\ell}\phi_{x\tau}\Big|^2
+ \Big|\frac{\epsilon}{\tilde\ell}\phi_{\tau}\Big|^2 + \Big|\frac{\epsilon^2}{\tilde\ell^2}\phi_{\tau\tau}\Big|^2 \, \Bigg] \,{\mathrm d}x {\mathrm d}\tau \right)^{\frac12},
\end{align}
with ${\tilde\ell}$ given in \eqref{lambda*}.
Moreover,
we have
\begin{align} \label{betanorm h2 mathfrak S estimate}
 \|\beta\phi\|_{H^2_*(\mathfrak S)} \leq C \|\beta h\|_{L^2(\mathfrak S)}.
\end{align}

\end{lemma}

\begin{proof}
 We express $h$ and $\phi$ as
\begin{align} \label{fourier-h}
 h(x, \tau) = \sum_{k=0}^\infty \Big[h_{1k} (x)\cos k\tau + h_{2k} (x)\sin k\tau\Big],
 \\[2mm]
 \label{fourier-phi}\phi(x, \tau) = \sum_{k=0}^\infty \Big[\phi_{1k} (x)\cos k\tau + \phi_{2k} (x)\sin k\tau\Big].
\end{align}
 Then we can  get
\begin{align} \label{linear philk1}
 - \frac{\epsilon^2k^2 4\pi^2}{{\tilde\ell}^2}\phi_{jk}+ \phi_{jk, xx} - \phi_{jk}+3U^2\phi_{jk} = h_{jk}\quad x\in {\mathbb R},
\end{align}
and $\phi_{jk}$ satisfies the orthogonality conditions
\begin{align} \label{linear philk2}
 \int_{\mathbb R}\phi_{jk}{\mathcal Z} \,{\mathrm d}x= \int_{\mathbb R}\phi_{jk} U' \,{\mathrm d}x \,=\, 0,
\end{align}
where $j=1,2$ and $k=0,1,2,\cdots$.

We first derive the estimates for $\phi_{jk}$ with $j=1,2$ and $k=0,1,2,\cdots$.

\smallskip\noindent$\clubsuit$\
 Consider the bilinear form
\begin{align*}
 B(\psi, \psi) = \int_{\mathbb R}\Big[|\psi_x|^2+ (1-3U^2)|\psi|^2 \Big] \,{\mathrm d}x.
\end{align*}
Since \eqref{linear philk2}, we can get
\begin{align} \label{bilinear 1}
 \|\phi_{jk}\|^2_{L^2({\mathbb R})} + \|\phi_{jk, x}\|^2_{L^2({\mathbb R})} \leq C B(\phi_{jk}, \phi_{jk}).
\end{align}

\smallskip\noindent$\clubsuit$\
Using \eqref{linear philk1} and \eqref{bilinear 1}, we can  claim that
\begin{align} \label{lin7-12}
\Bigg[1+\frac{\epsilon^4k^4}{{\tilde\ell}^4}\Bigg] \|\phi_{jk}\|^2_{L^2({\mathbb R})}
+ \Bigg[1+\frac{\epsilon^2k^2}{{\tilde\ell}^2}\Bigg] \|\phi_{jk, x}\|^2_{L^2({\mathbb R})} \leq C\|h_{jk}\|^2_{L^2({\mathbb R})}.
\end{align}
In fact, from \eqref{linear philk1}, we have
\begin{align*}
 \Bigg[1+\frac{\epsilon^2k^2}{{\tilde\ell}^2}\Bigg] \|\phi_{jk}\|^2_{L^2({\mathbb R})} + \|\phi_{jk, x}\|^2_{L^2({\mathbb R})} \leq C\|h_{jk}\|^2_{L^2({\mathbb R})}.
\end{align*}
 If $ \frac{\epsilon^2k^2}{{\tilde\ell}^2} \leq 1$, then we have $ \frac{\epsilon^4k^4}{{\tilde\ell}^4} \leq \frac{\epsilon^2k^2}{{\tilde\ell}^2}$. We can obtain \eqref{lin7-12}  directly .

\noindent If $ \frac{\epsilon^2k^2}{{\tilde\ell}^2} \geq 1$, then from \eqref{linear philk1}, we have
\begin{align*}
 \frac{\epsilon^2k^2}{{\tilde\ell}^2} B(\phi_{jk}, \phi_{jk}) + \frac{\epsilon^4k^4}{{\tilde\ell}^4} \|\phi_{jk}\|^2_{L^2({\mathbb R})} \leq C\|h_{jk}\|^2_{L^2({\mathbb R})} + (1-c)\frac{\epsilon^4k^4}{{\tilde\ell}^4} \|\phi_{jk}\|^2_{L^2({\mathbb R})},
\end{align*}
 i.e.,
\begin{align*}
 \frac{\epsilon^2k^2}{{\tilde\ell}^2} B(\phi_{jk}, \phi_{jk}) + c\frac{\epsilon^4k^4}{{\tilde\ell}^4} \|\phi_{jk}\|^2_{L^2({\mathbb R})} \leq C\|h_{jk}\|^2_{L^2({\mathbb R})},
\end{align*}
with $c\in (0, 1)$.
 Then by \eqref{bilinear 1}, we get \eqref{lin7-12}.

\smallskip\noindent$\clubsuit$\
By \eqref{lin7-12} and \eqref{linear philk1}, we can directly get
\begin{align}\label{second order estiamte}
\|\phi_{jk, xx}\|^2_{L^2({\mathbb R})} \leq C\|h_{jk}\|^2_{L^2({\mathbb R})}.
\end{align}


 From the expression of $\phi$ as in \eqref{fourier-phi}, we have
\begin{align*}
 \phi_x(x, \tau) &= \sum_{k=0}^\infty \Big[\phi_{1k, x} (x)\cos k\tau + \phi_{2k, x} (x)\sin k\tau\Big],
 \\[2mm]
 \phi_{xx}(x, \tau) &= \sum_{k=0}^\infty \Big[\phi_{1k, xx} (x)\cos k\tau + \phi_{2k, xx} (x)\sin k\tau\Big],
 \\[2mm]
 \phi_\tau(x, \tau) & = k\sum_{k=0}^\infty \Big[-\phi_{1k} (x)\sin k\tau + \phi_{2k} (x)\cos k\tau\Big],
 \\[2mm]
 \phi_{\tau\tau}(x, \tau) & = -k^2\sum_{k=0}^\infty \Big[\phi_{1k} (x)\cos k\tau + \phi_{2k} (x)\sin k\tau\Big],
 \\[2mm]
 \phi_{x\tau}(x, \tau) & = k\sum_{k=0}^\infty \Big[-\phi_{1k, x} (x)\sin k\tau + \phi_{2k,x} (x)\cos k\tau\Big].
\end{align*}
Therefore, we can derive that
\begin{align*}
 \|\phi\|_{L^2(\mathfrak S)}^2 & =\pi \sum_{k=0}^\infty \left[\|\phi_{1k}\|^2_{L^2({\mathbb R})}+ \|\phi_{2k}\|^2_{L^2({\mathbb R})}\right],
 \\[2mm]
 \|\phi_x\|_{L^2(\mathfrak S)}^2 & = \pi\sum_{k=0}^\infty \left[\|\phi_{1k, x}\|^2_{L^2({\mathbb R})}+ \|\phi_{2k, x}\|^2_{L^2({\mathbb R})}\right],
 \\[2mm]
 \|\phi_{xx}\|_{L^2(\mathfrak S)}^2 &= \pi \sum_{k=0}^\infty \left[\|\phi_{1k, xx}\|^2_{L^2({\mathbb R})}+ \|\phi_{2k, xx}\|^2_{L^2({\mathbb R})}\right],
 \\[2mm]
 \|\phi_\tau\|_{L^2(\mathfrak S)}^2 &= \pi \sum_{k=0}^\infty \left[k^2\|\phi_{1k}\|^2_{L^2({\mathbb R})}+ k^2\|\phi_{2k}\|^2_{L^2({\mathbb R})}\right],
 \\[2mm]
 \|\phi_{\tau\tau}\|_{L^2(\mathfrak S)}^2 &=\pi \sum_{k=0}^\infty \left[k^4\|\phi_{1k}\|^2_{L^2({\mathbb R})}+ k^4\|\phi_{2k}\|^2_{L^2({\mathbb R})}\right],
 \\[2mm]
 \|\phi_{x\tau}\|_{L^2(\mathfrak S)}^2 &= \pi \sum_{k=0}^\infty \left[k^2\|\phi_{1k, x}\|^2_{L^2({\mathbb R})}+ k^2\|\phi_{2k,x}\|^2_{L^2({\mathbb R})}\right].
\end{align*}
 Then from \eqref{lin7-12}, \eqref{second order estiamte}, we have
\begin{align*}
 &\|\phi\|_{L^2(\mathfrak S)}^2 + \|\phi_x\|_{L^2(\mathfrak S)}^2 + \|\phi_{xx}\|_{L^2(\mathfrak S)}^2+ \frac{\epsilon^2}{\tilde\ell^2} \|\phi_\tau\|_{L^2(\mathfrak S)}^2 + \frac{\epsilon^4}{\tilde\ell^4} \|\phi_{\tau\tau}\|_{L^2(\mathfrak S)}^2 + \frac{\epsilon^2}{\tilde\ell^2} \|\phi_{x\tau}\|_{L^2(\mathfrak S)}^2
\nonumber\\[2mm]
 & \leq C\sum_{k=0}^\infty \|h_{jk}\|^2_{L^2({\mathbb R})} \leq  C\|h\|_{L^2(\mathfrak S)}^2.
\end{align*}
This finishes the proof of \eqref{norm h2 mathfrak S estimate}.


 Next, we consider the proof of \eqref{betanorm h2 mathfrak S estimate}.
By denoting \[ \hat\phi = \beta \phi, \] we have,
for $\,(x, \tau)\in \mathfrak S$,
\begin{align*}
 \frac{\epsilon^2 4\pi^2}{{\tilde\ell}^2} \Big(\frac{\hat\phi} {\beta} \Big)_{\tau\tau} + \frac1\beta\Big[\hat\phi_{xx} - \hat\phi+3U^2 \hat\phi \Big]= h,
%\quad \text{for}\,(x, \tau)\in \mathfrak S,
\end{align*}
i.e.,
\begin{align*}
\frac1\beta\Big[ \frac{\epsilon^2 4\pi^2}{{\tilde\ell}^2} \hat\phi_{\tau\tau} + \hat\phi_{xx} - \hat\phi+3U^2 \hat\phi \Big] + \frac{\epsilon^2 4\pi^2}{{\tilde\ell}^2}\Big[2\hat\phi_\tau \Big(\frac1\beta \Big)_\tau + \hat\phi \Big(\frac1\beta \Big)_{\tau\tau} \Big] = h,
\end{align*}
i.e.,
\begin{align*}
\frac1\beta\Big[ \frac{\epsilon^2 4\pi^2}{{\tilde\ell}^2} \hat\phi_{\tau\tau} + \hat\phi_{xx} - \hat\phi+3U^2 \hat\phi \Big] + \frac{\epsilon^2 4\pi^2}{{\tilde\ell}^2}\Big[2\hat\phi_\tau \Big(-\frac{\beta'}{\beta^2} \frac{\tilde \ell} {2\pi \beta}\Big) + \hat\phi \Big(-\frac{\beta''} {\beta^3} +3\frac{|\beta'|^2} {\beta^4} \Big)\frac{\tilde \ell^2} {4\pi^2\beta} \Big] = h,
\end{align*}
with
\begin{align*}
\Big(\frac1\beta \Big)_\tau = -\frac{\beta'}{\beta^2} \frac{\tilde \ell} {2\pi \beta},
\qquad
\Big(\frac1\beta \Big)_{\tau\tau} = \Big(-\frac{\beta''} {\beta^3} +3\frac{|\beta'|^2} {\beta^4} \Big)\frac{\tilde \ell^2} {4\pi^2\beta}.
\end{align*}
Then, by setting
\begin{align*}
 \tilde h
%  & = \beta h -\frac{\epsilon^2 4\pi^2}{{\tilde\ell}^2}\Big[2\hat\phi_\tau \big(-\frac{\beta'}{\beta^2} \frac{\tilde \ell} {2\pi}\big) + \hat\phi \big(-\frac{\beta''} {\beta^3} +3\frac{|\beta'|^2} {\beta^4} \big)\frac{\tilde \ell^2} {4\pi^2} \Big]
% \nonumber
% \\[2mm]
& =\beta h - \epsilon^2\Big[2\hat\phi_\tau \Big(-\frac{\beta'}{\beta^2} \frac{2\pi}{\tilde \ell} \Big) + \hat\phi \Big(-\frac{\beta''} {\beta^3} +3\frac{|\beta'|^2} {\beta^4} \Big) \Big],
\end{align*}
 we get
\begin{align*}
\frac{\epsilon^2 4\pi^2}{{\tilde\ell}^2} \hat\phi_{\tau\tau} + \hat\phi_{xx} - \hat\phi+3U^2 \hat\phi = \tilde h,
\end{align*}
with the same conditions in \eqref{linearphi2}-\eqref{linearphi3} by replacing $\phi$ with $\hat\phi$.
By \eqref{beta-bound}, we get \eqref{betanorm h2 mathfrak S estimate} directly from \eqref{norm h2 mathfrak S estimate}.
We finish the proof.
\end{proof}

\end{document}
